\newcommand{\figuraCafeRotante}{%
\begin{figure}[H]
    \centering
    \begin{tikzpicture}[>=Latex,font=\small,x=1.15cm,y=1.0cm,
                        line cap=round,line join=round]
        \draw[very thick,rounded corners=6pt]
            (0.8,6.6)--(0.8,0.6)--(6.8,0.6)--(6.8,6.6);

        \fill[pink!28] (0.88,0.68) rectangle (6.72,4.80);
        \draw[very thick,magenta!75!black] (0.88,4.80)--(6.72,4.80);
        \node[magenta!75!black,anchor=south] at (3.8,4.86)
            {superficie libre inicial};

        \draw[densely dashed] (3.8,0.6)--(3.8,6.9);
        \node[anchor=south] at (3.8,6.64) {eje de rotación};

        \draw[<->,thick] (0.18,0.6)--(0.18,4.80)
            node[midway,left] {$7\,\mathrm{cm}$};
        \draw[<->,thick] (0.18,4.80)--(0.18,6.6)
            node[midway,left] {$3\,\mathrm{cm}$};
        \draw[densely dashed] (0.3,4.80)--(0.88,4.80);
        \draw[densely dashed] (0.3,6.60)--(0.88,6.60);

        \draw[<->,thick] (3.8,0.05)--(6.8,0.05)
            node[midway,below] {$R=3\,\mathrm{cm}$};
        \draw[densely dashed] (3.8,0.2)--(3.8,0.6);
        \draw[densely dashed] (6.8,0.2)--(6.8,0.6);

        \node[align=center,text=black!65] at (3.8,-0.85)
            {profundidad inicial del café: $7\,\mathrm{cm}$\\
             espacio libre hasta el borde: $3\,\mathrm{cm}$};
    \end{tikzpicture}
    \caption{Estado inicial de la taza. La superficie del café es horizontal
    y se encuentra a $7\,\mathrm{cm}$ sobre el fondo. Como la taza tiene
    $10\,\mathrm{cm}$ de profundidad, quedan $3\,\mathrm{cm}$ hasta el borde.}
    \label{fig:cafe-inicial}
\end{figure}

\begin{figure}[H]
    \centering
    \begin{tikzpicture}[>=Latex,font=\small,x=1.15cm,y=1.0cm,
                        line cap=round,line join=round]
        \draw[very thick,rounded corners=6pt]
            (0.8,6.6)--(0.8,0.6)--(6.8,0.6)--(6.8,6.6);

        \fill[pink!32] (0.88,0.68)--(6.72,0.68)--
            plot[domain=6.72:0.88,samples=140]
                (\x,{3.0+0.40*(\x-3.8)*(\x-3.8)})--cycle;
        \draw[very thick,magenta!75!black,domain=0.88:6.72,samples=140]
            plot (\x,{3.0+0.40*(\x-3.8)*(\x-3.8)});

        \draw[densely dashed,black!65] (0.88,4.80)--(6.72,4.80);
        \node[anchor=south,text=black!65] at (3.8,4.84)
            {nivel inicial: $7\,\mathrm{cm}$};
        \draw[densely dashed] (3.8,0.6)--(3.8,6.95);
        \node[anchor=south] at (3.8,6.68) {eje de rotación};

        \draw[<->,red!75!black,thick] (3.15,3.0)--(3.15,4.80)
            node[midway,left] {descenso: $3\,\mathrm{cm}$};
        \node[red!75!black,anchor=east] at (3.08,3.0)
            {$z_c=4\,\mathrm{cm}$};
        \draw[densely dashed,red!75!black] (3.15,3.0)--(3.8,3.0);

        \draw[<->,blue!70!black,thick] (1.30,4.80)--(1.30,6.41)
            node[midway,left] {ascenso: $3\,\mathrm{cm}$};
        \node[blue!70!black,anchor=south west] at (1.36,6.41)
            {$z_b=10\,\mathrm{cm}$};

        \draw[<->,thick] (7.28,3.0)--(7.28,6.41)
            node[midway,right] {$h=6\,\mathrm{cm}$};
        \draw[densely dashed] (6.72,3.0)--(7.4,3.0);
        \draw[densely dashed] (6.72,6.41)--(7.4,6.41);

        \draw[<->,thick] (3.8,0.05)--(6.8,0.05)
            node[midway,below] {$R=3\,\mathrm{cm}$};
        \draw[densely dashed] (3.8,0.2)--(3.8,0.6);
        \draw[densely dashed] (6.8,0.2)--(6.8,0.6);

        \draw[->,ultra thick,purple!75!black] (3.8,1.25)
            arc[start angle=90,end angle=-235,radius=0.50];
        \node[purple!75!black] at (4.58,1.25) {$\Omega$};

        \fill[magenta!75!black] (1.15,0.95) circle (0.09);
        \node[magenta!75!black,anchor=west] at (1.28,0.95) {$A$};

        \node[align=center,text=black!65] at (3.8,-0.90)
            {el borde sube $3\,\mathrm{cm}$ y el centro baja $3\,\mathrm{cm}$;\\
             la altura media permanece en $7\,\mathrm{cm}$};
    \end{tikzpicture}
    \caption{Taza en la condición de derrame. La conservación del volumen
    exige que el borde ascienda desde $7$ hasta $10\,\mathrm{cm}$ y que el
    centro descienda desde $7$ hasta $4\,\mathrm{cm}$. La diferencia entre
    ambos niveles es $h=6\,\mathrm{cm}$.}
    \label{fig:cafe-rotante}
\end{figure}%
}